\section*{Lời nói đầu}
\thispagestyle{empty}
Trong thực tế, đặc biệt là trong lĩnh vực y tế, việc quản lý và tổ chức hàng đợi bệnh nhân một cách khoa học và hiệu quả là yếu tố then chốt giúp tiết kiệm thời gian, nâng cao chất lượng phục vụ và đảm bảo công bằng trong việc khám chữa bệnh. Tuy nhiên, nếu không có một hệ thống hợp lý, những bệnh nhân cần được ưu tiên có thể phải chờ đợi lâu hoặc bị lẫn lộn trong số đông, gây ra hậu quả đáng tiếc.

Xuất phát từ nhu cầu đó, chúng em đã lựa chọn và thực hiện đề tài \textbf{ “Xây dựng hệ thống quản lý hàng đợi khám bệnh theo mức độ ưu tiên”}. Mục tiêu chính của đề tài là thiết kế và xây dựng một hệ thống hỗ trợ quản lý bệnh nhân theo thứ tự ưu tiên như cấp cứu, rất khẩn cấp, thông thường, tái khám… đảm bảo rằng những bệnh nhân có tình trạng nghiêm trọng hơn hoặc đến sớm hơn sẽ được ưu tiên khám trước.

Trong quá trình thực hiện, nhóm chúng em đã áp dụng những kiến thức đã học trong môn \textbf{Cấu trúc dữ liệu và Giải thuật}, trong đó có các phần là:
\begin{itemize}
    \item Thiết kế kiểu dữ liệu trừu tượng (ADT) để biểu diễn các đối tượng như bệnh nhân và hàng đợi ưu tiên.
    \item Cài đặt Max-Heap để xử lý hàng đợi theo mức độ ưu tiên và thời gian đăng ký với hiệu năng tốt.
    \item Kết hợp thuật toán so sánh ưu tiên kép (priority và thời gian đến) để chọn bệnh nhân phù hợp nhất.
    \item Triển khai cơ chế chống chờ vô hạn để đảm bảo người có mức ưu tiên thấp không bị "bỏ quên" trong hệ thống.
\end{itemize}

Qua quá trình xác định và thực hiện đề tài này không chỉ giúp chúng em rèn luyện kỹ năng lập trình, tư duy thiết kế hệ thống, mà còn giúp hiểu rõ hơn về giá trị thực tiễn của các cấu trúc dữ liệu và thuật toán trong việc giải quyết những vấn đề phức tạp trong đời sống.

Hy vọng báo cáo này của chúng em sẽ thể hiện được đầy đủ quá trình tìm hiểu, xây dựng và vận dụng kiến thức để hoàn thành bài toán một cách hiệu quả.
\newpage